\begin{statementbox}
	\textbf{\textcolor{CStatement}{条件}}
	\begin{description}[style=nextline, leftmargin=2.8em, labelsep=0.5em, itemsep=0.35em]
		\item[\textbf{\textcolor{CStatement}{条件 1（闭区间单调）：}}]
			函数 $y=f(x)$ 在闭区间 $[a,b]$ 上有定义且单调。
	\end{description}
	\textbf{\textcolor{CStatement}{结论}}
	\begin{description}[style=nextline, leftmargin=2.8em, labelsep=0.5em, itemsep=0.45em]
		\item[\textbf{\textcolor{CStatement}{结论一（递增求值域）：}}]
			\textbf{\textcolor{CStatement}{直观描述：}}
			最小值在左端点、最大值在右端点。
			\[
				f([a,b]) = [f(a), f(b)]
			\]
		\item[\textbf{\textcolor{CStatement}{结论二（递减求值域）：}}]
			\textbf{\textcolor{CStatement}{直观描述：}}
			最大值在左端点、最小值在右端点。
			\[
				f([a,b]) = [f(b), f(a)]
			\]
	\end{description}
	\textbf{\textcolor{CStatement}{推广结论（分段单调求值域）：}}
	\begin{description}[style=nextline, leftmargin=2.8em, labelsep=0.5em, itemsep=0.45em]
		\item[\textbf{\textcolor{CStatement}{直观描述：}}]
			先找单调性改变点，分区间求值域后合并。
		\item[\textbf{\textcolor{CStatement}{适用范围：}}]
			适用于在区间内不整体单调但可分段单调的函数，如二次函数、含绝对值函数。
	\end{description}

\end{statementbox}
